% vim: ai sw=2 spell spelllang=cs

\begin{frame}{Dynamická alokace}{Motivace}

  \begin{itemize}
    \item U statické alokace jsou velikost a počet známy v době překladu
    \begin{itemize}
      \item \code{int array[100];}
      \item \code{int a, b, c;}
    \end{itemize}

    \pause

    \item Ne vždy jsou tyto informace známy
    \begin{itemize}
      \item Závislost na vstupu
      \item Např.\ řetězec má typicky 20 znaků, ale může mít i 1000
    \end{itemize}

    \pause

    \item Dynamická alokace
    \begin{itemize}
      \item Umožňuje vytvořit datovou položku až za běhu

      \pause

      \item Ale může selhat
      \begin{itemize}
	\item Linux je optimista -- selže málokdy
	\item \cmd{sysctl vm.overcommit_memory}
      \end{itemize}
    \end{itemize}
  \end{itemize}

\end{frame}

\begin{frame}{Statická vs.\ dynamická alokace}

  \begin{itemize}
    \item \emph{Statická} alokace alokuje na zásobníku
    \begin{itemize}
      \item Automaticky se uvolňuje po dokončení bloku kódu
      \begin{itemize}
	\item Typicky konec funkce, konec podmínky \code{if}, cyklu \code{for},
	  \dots
      \end{itemize}

      \pause

      \item Výrazně rychlejší (obsah zásobníku je typicky v cache)
      \item Využití pro krátkodobé proměnné/objekty

      \pause

      \item Např.\ \code{int array[100];}
    \end{itemize}

    \pause

    \item \emph{Dynamická} alokace na haldě
    \begin{itemize}
      \item Zůstává do explicitního uvolnění (nebo konce aplikace)
      \item Využití pro dlouhodobé paměťové entity

      \pause

      \item Např.\ \code{int *array = malloc(variable_len * sizeof(int));}
    \end{itemize}

    \pause

    \item Specialitou je \emph{alokace polí} s variabilní délkou (od C99)
    \begin{itemize}
      \item Délka není známa v době překladu
      \item Ale alokuje se na zásobníku a automaticky odstraňuje

      \pause

      \item Např.\ \code{int array[length];}
    \end{itemize}
  \end{itemize}

\end{frame}

\begin{frame}{Alokace dynamické paměti}

  \begin{itemize}
    \item Hlavička: \header{stdlib.h}

    \item \code{void *malloc(size_t size)}
    \begin{itemize}
      \item \code{size}: velikost k alokaci
      \item Návratová hodnota: NULL při chybě, jinak adresa

      \pause

      \item Jeden souvislý blok
      \item Na haldě
      \item Paměť není inicializována (,,smetí'')
    \end{itemize}

    \pause

    \item \code{void *calloc(size_t nmemb, size_t size)}
    \begin{itemize}
      \item \code{nmemb}: počet položek
      \item \code{size}: velikost položky
      \item Alokuje \code{nmemb * size}

      \pause

      \item Inicializuje paměť na 0
    \end{itemize}
  \end{itemize}

\end{frame}

\begin{frame}{Uvolnění dynamické paměti}

  \begin{itemize}
    \item \code{void free(void *ptr)}
    \begin{itemize}
      \item Uvolní alokovanou paměť
      \begin{itemize}
	\item \code{ptr} může být jen to, co vrátil \code{malloc} (i
	  \code{NULL})
	\item Tj.\ nelze uvolnit jen část \code{free(mem + 4)}
      \end{itemize}
    \end{itemize}

    \pause

    \item \pozor
    \begin{itemize}
      \item Nelze přistupovat na paměť po zavolání \code{free}
      \item \code{free} nemaže \emph{obsah} paměti
    \end{itemize}
  \end{itemize}

\end{frame}

\begin{frame}{Úkol}

  \heading{Alokace paměti}
  \begin{enumerate}
    \item Alokujte 1 GiB paměti
    \item Spěte 10 vteřin
    \begin{itemize}
      \item Mezitím pozorujte výstup \cmd{htop}/\cmd{top}/\cmd{free}
      \item Změnila se velikost volné paměti?
    \end{itemize}

    \item Vyzkoušejte
    \item Vysvětlete
  \end{enumerate}

\end{frame}

\begin{frame}{Změna dynamické paměti}

  \begin{itemize}
    \item \code{void *realloc(void *ptr, size_t size)}
    \begin{itemize}
      \item \code{(ptr == NULL)} $\Rightarrow$ \code{malloc(size)}
      \item \code{(size == 0)} $\Rightarrow$ \code{free(ptr)}
      \item \code{(ptr != NULL)} $\Rightarrow$ změna velikosti dynamické paměti
	na \code{size}

      \pause

      \item Návratová hodnota: nemusí být stejná jako \code{ptr}
    \end{itemize}

    \pause

    \item Obsah původního paměťového bloku je zachován

    \pause

    \item Při zvětšení je dodatečná paměť neinicializovaná
    \begin{itemize}
      \item Stejně jako při \code{malloc}
    \end{itemize}
  \end{itemize}

\end{frame}

\begin{frame}[fragile]{\code{memset}}

  \begin{itemize}
    \item Nastaví paměť na zadanou hodnotu

    \pause

    \item \code{void *memset(void *s, int c, size_t n)} (\header{string.h})
    \begin{itemize}
      \item \code{s} -- paměť k inicializaci
      \item \code{c} -- hodnota (znak), kterou inicializovat
      \item \code{n} -- velikost bloku k inicializaci
      \begin{itemize}
	\item Často \code{sizeof()}
      \end{itemize}
    \end{itemize}

    \pause

    \item Řádově rychlejší než \code{for} cyklus
  \end{itemize}

  \pause

  \begin{block}{Příklad}
  \begin{lstlisting}
struct {
  int a, b;
} *array = malloc(sizeof(*array));
memset(array, 0, sizeof(*array));
  \end{lstlisting}
  \end{block}

\end{frame}

\begin{frame}{Dynamická alokace}{Shrnutí}

  \begin{enumerate}
    \item Alokujeme potřebný počet bajtů
    \begin{itemize}
      \item Pozor na správnou velikost
      \item Typicky používáme \code{pocet_prvku * sizeof(prvek)}
    \end{itemize}

    \pause

    \item Uložíme adresu paměti do ukazatele s požadovaným typem
    \begin{itemize}
      \item \code{int *array = malloc(100 * sizeof(*array))}
    \end{itemize}

    \pause

    \item Po konci práce uvolníme
    \begin{itemize}
      \item \code{free(array)}
    \end{itemize}
  \end{enumerate}

\end{frame}

